\chapter{ESTADO DEL ARTE Y REVISIÓN CRÍTICA DE LA LITERATURA}

\section{Panorama general}
El control adaptativo de semáforos con Aprendizaje por Refuerzo (RL) se ha consolidado como una alternativa a la semaforización fija y actuada, al permitir políticas que se adaptan en tiempo real a la variación de flujos. La literatura reciente estudia tanto agentes únicos por intersección (single-agent) como enfoques multiagente (un agente por intersección) y evalúa su impacto en métricas como retraso promedio, colas, throughput y emisiones. En esta sección se sintetizan tres trabajos directamente relevantes a nuestro proyecto (Uniandes 2023, Trinity 2022, UAQ 2024) y se discuten sus aportes, limitaciones y la pertinencia para Lima Metropolitana.

\section{Estado del arte: trabajos estrechamente relacionados}

\subsection*{Universidad de los Andes (2023): \textit{Semáforos inteligentes por medio de aprendizaje por refuerzo} \citep{Chavez2023}}
\textbf{Objetivo y configuración.} Estudiar el desempeño de algoritmos RL (Q-learning, SARSA, DQN y PPO) para reducir tiempos de espera frente a un control de tiempos fijos en un escenario simulado. \\
\textbf{Métodos.} Implementación de algoritmos tabulares (Q-learning, SARSA) y de Deep RL (DQN, PPO) con política $\varepsilon$-greedy en tabulares y objetivos recortados en PPO; comparación contra baseline de tiempos fijos. \\
\textbf{Métricas.} Eficiencia relativa frente a tiempos fijos (reducción del tiempo de espera), número de vehículos en sistema, estabilidad de la política. \\
\textbf{Resultados.} Reporta eficiencias de \textbf{35.79\% (Q-learning), 51.63\% (SARSA), 63.40\% (PPO) y 63.49\% (DQN)} frente a control fijo, medidas como reducción del \emph{tiempo de espera} y no del retraso. La distinción importa al comparar con esta tesis, que fija el retraso como métrica primaria y muestra que ambas magnitudes no varían en la misma proporción. DQN y PPO muestran mejor adaptación y menor tiempo de espera, con DQN muy ligeramente superior en el caso analizado. \\
\textbf{Limitaciones.} Demanda y geometría sintetizadas; no integra datos reales de conteo/OD ni perturbaciones estocásticas complejas; sin análisis estadístico robusto (p.\,ej., ICs o pruebas de significancia). \\
\textbf{Lecciones para esta tesis.} (i) Incluir \textit{deep baselines} (DQN/PPO) junto a tabulares; (ii) reportar mejoras con métricas comparables; (iii) añadir pruebas de significancia y sensibilidad.

\subsection*{Trinity College Dublin (2022): \textit{Reinforcement Learning for Traffic Light Optimization} \citep{Chakrabarti2022}}
\textbf{Objetivo y configuración.} Comparar agentes fijos/aleatorios con \textbf{DQN}, \textbf{A2C} y \textbf{PPO} en casos single-agent, multi-agent y un caso más realista. \\
\textbf{Métodos.} SUMO+TraCI, definición de estados (colas/velocidades/ocupación), espacios de acción (fases/tiempos) y recompensas (tiempos de espera acumulados/negativos). \\
\textbf{Métricas.} \textbf{Tiempo de espera acumulado} por intersección vs. pasos de simulación; comparación entre agentes y escenarios. \\
\textbf{Resultados.} \textbf{PPO} es consistentemente el mejor, seguido de \textbf{A2C}, agentes fijos/aleatorios; \textbf{DQN} queda por detrás en sus pruebas. \\
\textbf{Limitaciones.} Resultados centrados en comparativa algorítmica; no se usan datos reales de OD, ni se reportan pruebas de significancia. \\
\textbf{Lecciones para esta tesis.} (i) PPO y A2C son fuertes candidatos para entornos urbanos; (ii) conviene incorporar escenarios multiagente; (iii) sumar análisis estadístico y casos con datos reales para mayor validez externa.

\subsection*{UAQ (2024): \textit{Predicción de tráfico vehicular aplicada a la optimización semafórica en un microsimulador} \citep{UAQ2024}}
\textbf{Objetivo y configuración.} Integrar \textbf{predicción de demanda} (AR, Filtros de Kalman, RNN) con \textbf{microsimulación} para \textbf{optimizar ciclos semafóricos} en un corredor real. \\
\textbf{Datos.} Históricos de \textbf{7 intersecciones} (ago–sep 2022), con escenarios de 1 hora, 1 día, 1 semana y 1 mes. \\
\textbf{Métricas.} Errores de predicción (RMSE, MAE, MAPE, $R^2$) y, tras aplicar los ciclos en el simulador, tiempo de recorrido y colas. \\
\textbf{Resultados.} \textbf{RNN y Filtros de Kalman} alcanzan \textbf{$R^2\approx 0.95$} en horizontes corto/medio; al aplicar los ciclos óptimos, se obtiene \textbf{hasta 36.34\% de reducción del tiempo de viaje}. \\
\textbf{Limitaciones.} Optimización semafórica basada en pronóstico (no en RL); generalización a otras franjas/variabilidad no evaluada; requiere calidad de datos. \\
\textbf{Lecciones para esta tesis.} (i) Incorporar \textbf{datos reales (OD, aforos)} y \textbf{predicción de demanda} como \textit{módulo upstream} que alimente al agente RL; (ii) evaluar impacto sobre \textbf{tiempos de viaje} además del retraso.

\section{Revisión crítica y vacíos identificados}
\textbf{Modelado de la realidad.} Uniandes y Trinity validan con escenarios simulados; UAQ sí usa datos reales, pero sin RL. Esta tesis busca combinar ambas fortalezas, integrando datos reales de aforos sobre un cruce de Lima; en la fase que reporta este documento la validación es todavía sobre escenarios sintéticos. \\
\textbf{Línea base de comparación.} Los tres trabajos miden únicamente frente a control de tiempo fijo (Trinity añade un agente aleatorio). Ninguno compara contra semaforización \textit{actuada}, que es el control adaptativo estándar y está disponible de fábrica en SUMO. Esta tesis lo incorpora como segundo control de referencia, con las mismas restricciones de verde que el agente, y el Capítulo~\ref{cap:resultados}muestra que la omisión no es menor: el actuado supera al agente tabular en los cuatro escenarios evaluados. \\
\textbf{Comparación algorítmica.} Trinity muestra ventajas de PPO y A2C; Uniandes añade DQN con buen rendimiento. Se implementarán \textbf{Q-learning, SARSA, DQN y PPO}, con los tabulares como línea base y los profundos como comparación principal. \\
\textbf{Rigor experimental.} Falta análisis de significancia y de sensibilidad. La evaluación de esta tesis incluye \textbf{intervalos de confianza al 95\%}, prueba de Wilcoxon pareada por semilla, Mann--Whitney y tamaño de efecto, con semillas de evaluación disjuntas de las de entrenamiento. \\
\textbf{Transferencia a Lima.} Se escogerá una intersección o corredor representativo de Lima e integrará \textbf{restricciones operativas} (tiempos mínimos de verde y ámbar, coordinación en arterias) y un agente por intersección.

\section{Síntesis comparativa}
\begin{table}[H]
\centering
\footnotesize
\setlength{\tabcolsep}{6pt}
\renewcommand{\arraystretch}{1.2}
\begin{tabularx}{\textwidth}{@{} l c c X X @{}}
\toprule
\textbf{Trabajo} & \textbf{Datos} & \textbf{Simulador} & \textbf{Algoritmos} & \textbf{Resultados clave} \\
\midrule
\makecell[l]{Uniandes (2023)\\ \citep{Chavez2023}}
  & Sintéticos
  & SUMO
  & Q-learning, SARSA, DQN, PPO
  & Eficiencias vs. fijo: 35.79\%, 51.63\%, 63.40\%, 63.49\% \\
\makecell[l]{Trinity (2022)\\ \citep{Chakrabarti2022}}
  & Sintéticos
  & SUMO
  & DQN, A2C, PPO
  & PPO $>$ A2C $>$ (fijo/aleatorio) $>$ DQN en tiempo de espera \\
\makecell[l]{UAQ (2024)\\ \citep{UAQ2024}}
  & \textbf{Reales} (7 inters.)
  & SUMO
  & AR, Kalman, RNN + optimización de ciclos
  & $R^2\approx 0.95$ (corto/medio); \textbf{-36.34\%} en tiempo de viaje \\
\bottomrule
\end{tabularx}
\caption{Comparación de trabajos relacionados con control/optimización semafórica.}
\label{tab:relacionados}
\end{table}


\section{Implicancias para el diseño de esta tesis}

A partir del estado del arte, esta tesis implementa: (i) un entorno \textbf{SUMO+TraCI} con un agente por intersección, (ii) \textbf{Q-learning y SARSA} como líneas base tabulares y \textbf{DQN y PPO} como algoritmos de \textit{Deep RL}, y (iii) evaluación rigurosa (retraso, colas, \textit{throughput}, paradas y emisiones) frente a control fijo y actuado, con pruebas estadísticas. En la fase reportada en este documento está implementado Q-learning tabular; el resto queda como trabajo siguiente.

\noindent\textbf{Aporte específico (U4: Coordinación distribuida).} Se propone un módulo de \textbf{coordinación distribuida entre semáforos} para pasar de optimizaciones locales a \textbf{optimización de corredor}. El módulo está formulado y no se ejecuta en esta fase; los escenarios de varios semáforos evaluados en el Capítulo~\ref{cap:resultados}usan agentes independientes, y el hecho de que pierdan frente a un programa fijo coordinado es la evidencia empírica que lo motiva. Sus componentes previstos son:
\begin{itemize}
    \item \textit{Estado extendido}: cada agente observa su entorno local (colas por carril, fase, tiempo en fase) y un \textbf{resumen de vecinos} (fase vecina, tiempo restante, estimación de \emph{outflow} hacia su enlace compartido).
    \item \textit{Mensajería ligera}: canal de comunicación simulado (latencia/pérdida configurables) para intercambiar mensajes agent-to-agent (\emph{fase}, \emph{outflow esperado}, \emph{colas en salidas}).
    \item \textit{Recompensa con coordinación}: término explícito de \textbf{spillback} para evitar que una cola bloquee a la intersección vecina; opcionalmente mezcla con una señal global (\(R_{global}\)).
    \item \textit{Algoritmos}: \textbf{IPPO} (PPO independiente con pesos compartidos) como base; extensión a \textbf{MAPPO} (crítico centralizado) si el tiempo lo permite.
    \item \textit{Evaluación a nivel corredor}: tiempo de viaje extremo a extremo, \textit{throughput} total, longitud de cola máxima y tasa de \textit{spillback}, además de métricas locales.
\end{itemize}

Este módulo añade el tópico de \textbf{Sistemas Distribuidos} al marco de RL, incrementa el esfuerzo computacional (entrenamiento multiagente con estado ampliado) y permite medir el beneficio de la \textbf{coordinación} frente a agentes aislados, que es la pregunta que dejan abierta los resultados del Capítulo~\ref{cap:resultados}.
